\documentclass[10pt]{article}
\usepackage[margin=0.9in]{geometry}
\usepackage{url}
\setlength{\parindent}{0pt}\setlength{\parskip}{0.6em}
\begin{document}
October 2026

Editors, \emph{INFORMS Journal on Computing}

Dear Editors,

Please consider the manuscript \textbf{``Certifying rare overflow risk of regional pickup-and-delivery plans''} for publication in the \emph{INFORMS Journal on Computing}.

\textbf{What the paper does.} Distribution of medical, cold-chain and dangerous goods often requires that a routing plan overflow its vehicles with probability at most a service level between $10^{-3}$ and $10^{-6}$. We study how to verify such a level for a given plan of the vehicle routing problem with simultaneous pickup and delivery, under gamma demands with a Gaussian factor copula. We show that a route overflows through one of two end modes, that a single importance-sampling shift misses the weaker one and underestimates by about $0.8$ times its computed mass, and that a mixture at the two computed dominating points removes the bias without pilot simulation. The method extends to route days with stochastic travel times and to hubs whose routes share demand factors. It is checked on 148 instances (6,242 routes) from Northern Vietnam and Guangdong against a quadrature reference, crude Monte Carlo and budget-fair classical rivals (cross-entropy mixtures, Glasserman--Li, subset simulation), including coverage studies at $10^{-5}$ and $10^{-6}$. The plans we certify, built for a per-route risk of $0.10$, do not meet any level at their planned capacity and need $35$--$76\%$ more capacity. A short section tests quantum amplitude estimation on the same problem and reports that it does not beat the classical estimator here.

\textbf{Why IJOC.} The contribution is a computational method with an analysis of when it works, a reproducible implementation and a large experimental study, which is the combination the journal asks for. The paper sits between simulation methodology (rare-event estimation, variance reduction, confidence intervals) and its use in operations research (routing plans, chance constraints). The journal has published work in this line, for example T.\ Homem-de-Mello, ``A study on the cross-entropy method for rare-event probability estimation'', \emph{INFORMS Journal on Computing} 19(3):381--394 (2007), which studies importance sampling for rare-event probabilities in static simulation models. All code, instances, plans and the scripts behind every table and figure are public at \url{https://github.com/vinhqdang/stochastic_vrp}, and we will release them under the IJOC software policy.

\textbf{Area.} We submit to the area \emph{Simulation, Stochastic Models, \& Stochastic Optimization}: the core of the paper is importance-sampling estimation, with certified intervals, of a rare event of a stochastic model.

\textbf{Related manuscripts and overlap.} The submission uses instances and plan files that also appear in two of our other manuscripts, one under review at \emph{Computers \& Operations Research} (online mid-route recourse) and one under review at \emph{Transportation Science} (e-process monitoring). Their results are different: this paper contains no recourse policy and no e-values, and reports no number taken from them. A separate manuscript on quantum amplitude estimation for stochastic routing is in preparation; the quantum section here is a short companion check with a negative conclusion. This manuscript has not been submitted elsewhere.

Sincerely,\\[0.3em]
Quang-Vinh Dang\\
British University Vietnam, Hung Yen, Vietnam\\
vinh.dq4@buv.edu.vn, ORCID 0000-0002-3877-8024
\end{document}
